\subsection{First momentum correction to the primitive mode sum}
\label{sec:lovelock-momentum-time}

We now track the \(k^2\) derivative at fixed \(\nu\) in the
large-overtone regime. Define the optical symbol and the
logarithmic reflection derivative by
\begin{equation}
 \xi=e^{2\ii\omega\mathcal Z_{\rm L}-2\pi\ii h},\qquad
 \ell_{\rm A}(\omega)=
 \frac{e^{-\ii\pi/3}\Gamma(1/3)^2}
 {4\,3^{1/3}(\omega/2)^{2/3}} .
 \label{eq:lovelock-momentum-symbol}
\end{equation}
The symbol equals one at each positive-real-part pole at
\(k=0\). The displaced frequency
\(\omega=\omega_n^0+(\theta-\ii\log\chi_0)/(2\mathcal Z_{\rm L})\)
gives \(\xi=\chi_0 e^{\ii\theta}\) for a fixed \(0<\chi_0<1\).
The derivative \(\ell_{\rm A}\) follows from
\eqref{eq:lovelock-perturbed-airy} and
\eqref{eq:lovelock-airy-moment}.

In the leading matched expansion, the near-boundary
ingoing and outgoing waves are proportional to
\(r^{-1}H_\nu^{(1,2)}(\sqrt{\omega^2-k^2}/r)\).
If the outgoing-to-ingoing coefficient in this basis is
\(\Xi\), their small-argument expansions give the
response ratio \((1-e^{2\pi\ii\nu}\Xi)/(1-\Xi)\) times
the vacuum value \cite{Valach:2026,Jia:2026}.
At \(k=0\), the gamma reflection formula applied to
\eqref{eq:detuned-g} gives this same ratio with \(\Xi=\xi\),
multiplied by a gamma envelope tending to its vacuum value.
To first order in \(k^2\), the core Stokes coefficient gives
\(\partial_{k^2}\log\Xi=\ell_{\rm A}+o(|\omega|^{-2/3})\).
Outside the core matching radius \(\varepsilon_z\),
the variation of the transport action is
\(O(|\omega|^{-1}\varepsilon_z^{-1/2})\).
With \(\varepsilon_z|\omega|^{2/3}\to\infty\) this is
subleading, and the near-boundary vacuum response varies
at relative order \(k^2/\omega^2\).

Consequently the leading Airy prediction on bounded
frequency displacements from \(\omega_n^0\), away from
the poles in the displacement variable, is
\begin{align}
 \frac{\dot{\mathcal G}(\omega)}
      {\mathcal K_\nu(\omega)\ell_{\rm A}(\omega)}
 &=\frac{\xi}{(1-\xi)^2}+o(1),\notag\\
 \mathcal K_\nu(\omega)
 &=\frac{2\pi\ii}{\Gamma(\nu)\Gamma(1+\nu)}
                      \left(\frac{\omega}{2}\right)^{2\nu}.
 \label{eq:lovelock-momentum-transfer}
\end{align}
Here \(\mathcal K_\nu\) combines the vacuum amplitude
and the factor \(1-e^{2\pi\ii\nu}\) from differentiating
the boundary ratio. The branch is continued from real
positive \(\omega\). On \(|\xi|<1\), the coefficient of
\(\xi^m\) in \eqref{eq:lovelock-momentum-transfer} is \(m\).
The extra factor is the derivative of the \(m\)-fold
reflection amplitude.

\begin{table}[t]
 \centering\small
 \begin{tabular}{ccccc}
 \hline
 \(\nu\) & \(n\) & \(\lvert c_2-1\rvert\) & \(\lvert c_3-1\rvert\) &
   \(\lvert\dot{\mathcal R}_n^+/(\mathcal R_n^+b_n)\rvert\) \\
 \hline
 0.75 & 8 & 1.620e-03 & 8.598e-04 & 2.398e-02 \\
0.75 & 16 & 7.111e-04 & 3.673e-04 & 1.347e-02 \\
0.75 & 32 & 2.985e-04 & 1.526e-04 & 7.298e-03 \\
1.25 & 8 & 2.381e-03 & 3.670e-03 & 5.594e-02 \\
1.25 & 16 & 1.084e-03 & 1.660e-03 & 3.124e-02 \\
1.25 & 32 & 4.635e-04 & 7.069e-04 & 1.679e-02 \\
\hline

 \end{tabular}
 \caption{Test of \eqref{eq:lovelock-momentum-transfer} using
 finite-momentum boundary matching. The coefficients \(c_m\)
 are the projected harmonics divided by \(m\).}
 \label{tab:lovelock-momentum-harmonics}
\end{table}

For table~\ref{tab:lovelock-momentum-harmonics}, we project
the left-hand side of \eqref{eq:lovelock-momentum-transfer}
on harmonics \(\xi^2\) and \(\xi^3\) with \(\chi_0=0.45\).
We extract the moving-pole residue derivative by Cauchy
quadrature around the unperturbed pole. To understand its
suppression, write \(u=2\ii\mathcal Z_{\rm L}(\omega-\omega_n^0)\).
Then
\begin{equation}
 \frac{e^u}{(1-e^u)^2}=\frac{1}{u^2}+O(1).
\end{equation}
The leading normalized response has a double pole with
zero simple-pole coefficient. The normalizing factor in
\eqref{eq:lovelock-momentum-transfer} varies by a relative
\(O(N^{-1})\) on a circle of fixed radius, so the leading
matched expansion gives
\(\dot{\mathcal R}_n^+=o(\mathcal R_n^+ b_n)\).
This is the relative suppression tested in the last column.

The polar part of the perturbed time response is obtained
by differentiating the mode sum. For the positive-real-part
branch the term is
\begin{equation}
 -\ii\sum_{n=0}^\infty
   (\dot{\mathcal R}_n^+-\ii t\mathcal R_n^+ b_n)
     e^{-2(1+\ii)Nt}.
 \label{eq:lovelock-time-momentum-polar}
\end{equation}
The negative-real-part branch has the corresponding conjugate
time dependence, and both sums converge locally uniformly
in \(\operatorname{Re}t>|\operatorname{Im}t|\).
Near the upper primitive time \(t_m=m\tau(1+\ii)\),
the pole-motion term controls the leading singularity.
Using \(\mathcal R_n^+\sim C_{\rm R}N^{2\nu}\) and
\(b_n\sim\mathcal B_{\rm A}N^{-2/3}\) in the rescaled sum gives
\begin{equation}
 \left.\dot{\mathcal G}_R(t)\right|_{\rm polar}
 \ \sim\
 -t_m C_{\rm R}\mathcal B_{\rm A} e^{-2\pi\ii m h}
 \frac{\Gamma(2\nu+1/3)}
      {[2(1+\ii)(t-t_m)]^{2\nu+1/3}} .
 \label{eq:lovelock-momentum-time-power}
\end{equation}
The singular power decreases by \(2/3\) from
\eqref{eq:lovelock-retarded-time}. The linear dependence
on \(t_m\) carries the factor of \(m\) in the
frequency-space harmonics. In fact \(t_m=2m\mathcal Z_{\rm L}\),
and \(\mathcal B_{\rm A}=\ii\ell_{\rm A}(2(1-\ii))/(2\mathcal Z_{\rm L})\).
Dividing by \eqref{eq:lovelock-retarded-time} gives
\begin{equation}
 \frac{\left.\dot{\mathcal G}_R(t)\right|_{\rm polar}}
      {\mathcal G_R(t)}
 \sim
 \frac{m e^{-\ii\pi/3}\Gamma(1/3)^2}{4\,3^{1/3}}
 \frac{\Gamma(2\nu+1/3)}{\Gamma(2\nu+1)}
 [2\ii(t-t_m)]^{2/3}.
 \label{eq:lovelock-momentum-time-ratio}
\end{equation}
This is the same relative coefficient obtained by Laplace
transforming the factor \(m\ell_{\rm A}(\omega)\) in the
isolated \(m\)-reflection term.

\begin{table}[t]
 \centering\small
 \begin{tabular}{cccc}
 \hline
 \(m\) & \(\epsilon=0.15\) & \(\epsilon=0.075\) & \(\epsilon=0.035\) \\
 \hline
 2 & 3.941e-02 & 1.758e-02 & 7.734e-03 \\
3 & 2.926e-02 & 1.228e-02 & 5.185e-03 \\
\hline

 \end{tabular}
 \caption{Relative error of the leading coefficient in
 \eqref{eq:lovelock-momentum-time-power}, using the exact
 pole-motion terms at \(\nu=0.75\) and \(t=t_m+\epsilon\).}
 \label{tab:lovelock-momentum-time}
\end{table}

Table~\ref{tab:lovelock-momentum-time} uses the first 320
exact pole slopes and zero-momentum residues to test the
coefficient at the second and third primitive times.
Equation \eqref{eq:lovelock-momentum-time-power} concerns the
polar contribution at order \(k^2\).
Identifying it with the entire Fourier continuation requires
growth estimates for the response derivative away from the poles.
